% ============================================================
% Section 05 — Conditions of Identity Preservation
% Book I: Ontology of Viability
% Chapter 05: Identity and Continuity
%
% Purpose:
%   To investigate the ontological conditions under which
%   the preservation of identity through transformation
%   may be established.
%
% Conceptual Policy:
%   1. Identity preservation must not be reduced to a
%      single material, structural, functional, relational,
%      or invariant condition.
%   2. No proposed condition is universally necessary
%      or sufficient without independent justification.
%   3. Identity preservation must be distinguished from
%      persistence, similarity, and continuity of manifestation.
%   4. Conditions of identity preservation must be examined
%      relative to the existent, domain, and respect involved.
%   5. The possibility of transformation without identity
%      loss must remain distinct from an unrestricted
%      assumption of identity persistence.
% ============================================================

\section{Conditions of Identity Preservation}
\label{sec:conditions-of-identity-preservation}

\subsection{The Ontological Problem of Identity Preservation}

The question of identity preservation concerns the
conditions under which an existent remains numerically
the same existent through transformation.

The preceding sections have established that identity
cannot be reduced to persistence, invariance, or the
continuity of particular manifestations. They have
also shown that qualitative difference does not
automatically establish numerical non-identity.

These distinctions, however, do not yet provide
an account of the conditions under which identity
is preserved.

The problem requires an examination of the ontological
conditions that may justify the attribution of one
and the same identity across different conditions
of an existent.

Such conditions cannot be selected merely because
they are observable, relatively stable, or frequently
associated with continuity. Their relevance must
be established in relation to the identity under
examination.

Material constitution, structural organization,
determination, relational configuration, invariance,
and manifestation may all contribute to an account
of identity preservation. Nevertheless, none can
be treated in advance as a universal and sufficient
criterion.

The present section therefore investigates identity
preservation as an ontological problem whose conditions
must be established through analysis rather than
assumed through a preferred model of continuity.

\subsection{The Difference Between Identity Preservation and Feature Preservation}

Identity preservation must be distinguished from
the preservation of particular features.

An existent may preserve a feature while undergoing
transformations in other respects. It may also lose
or transform a previously present feature without
this fact alone establishing the termination of
its identity.

Feature preservation concerns the continuation
of a determination, relation, structure, or other
specified characteristic. Identity preservation
concerns the continued numerical identity of the
existent itself.

The two may be related, but their relation requires
independent justification.

A feature may be shared by numerically distinct
existents. Consequently, its preservation cannot
by itself establish that the existent exhibiting
it remains numerically identical with an earlier
existent.

Likewise, the transformation of a feature does
not independently establish that the existent
has become numerically different.

The investigation must therefore distinguish between
features that remain preserved and the ontological
conditions under which the identity of their bearer
is preserved.

\subsection{Identity Preservation and the Ontological Unity of the Existent}

An account of identity preservation must examine
the relation between identity and the existent
whose identity is being investigated.

Identity is not an additional property that is
attached to an otherwise independently constituted
existent. Nor is the existent merely an aggregate
of features whose preservation automatically
establishes the preservation of identity.

The ontological question concerns the continued
identity of the existent through the transformations
that affect its determinations and conditions.

This requires distinguishing the unity of identity
from the unity of a configuration, structure,
material aggregate, or manifestation.

An existent may undergo changes in the organization
of its components or relations without these changes
alone determining whether its identity has been preserved.

Conversely, the apparent unity of a configuration
does not independently establish the numerical
identity of the existent constituted or manifested
through it.

The relation between identity and the unity of
the existent must therefore be examined without
reducing one to the other.

The conditions of identity preservation must account
for the existent as an ontological subject rather
than merely identifying a collection of features
that remain unchanged.

\subsection{Material Conditions of Identity Preservation}

Material constitution may be relevant to the
preservation of identity in particular domains.

An existent may depend upon a specific material
constitution for its continued existence or for
the realization of certain determinations. In such
cases, material conditions may contribute to the
conditions under which identity preservation is
examined.

However, the relevance of material constitution
does not establish that identity is universally
identical with material continuity.

Material constituents may be replaced, reorganized,
or transformed while the question of identity
remains open. Conversely, the preservation of
the same material constituents does not independently
establish numerical identity.

The material conditions of identity preservation
must therefore be examined in relation to the
kind of existent and the ontological role played
by its material constitution.

A material condition may be necessary in a particular
domain without being universally necessary. It
may contribute to identity preservation without
being sufficient to establish it.

The ontological status of material conditions must
be determined through a specific account of the
existent under examination.

\subsection{Structural Conditions of Identity Preservation}

Structural organization may also contribute to
the preservation of identity.

An existent may depend upon particular relations
among its components, determinations, or internal
configurations. The alteration of these relations
may affect the conditions under which its identity
is maintained.

Nevertheless, structural preservation cannot
be treated as a universal criterion of identity.

A structure may be reproduced in a numerically
distinct existent. Structural equivalence therefore
does not independently establish numerical identity.

Likewise, structural transformation does not
automatically establish the termination of identity.
An existent may undergo reorganization while the
ontological question of identity preservation
remains unresolved.

The relevant question is not merely whether a
structure has remained unchanged, but whether
the structure in question has an ontologically
constitutive role in the identity of the existent.

Such a role must be demonstrated for the particular
domain under examination.

\subsection{Determination and Identity Preservation}

Determinations may contribute to the conditions
under which identity preservation is investigated.

An existent is not without determination, and
transformations in its determinations may alter
the ways in which its existence is constituted
or manifested.

However, the preservation of a particular determination
does not independently establish identity preservation.
Nor does the alteration of a determination automatically
establish non-identity.

Determinations may differ in their ontological
significance. Some may be constitutive under
particular conditions, while others may be contingent
or transformable without determining the numerical
identity of the existent.

The investigation must therefore distinguish
between the preservation of determinations and
the preservation of identity.

It must also avoid assuming that every determination
has the same relation to identity or that identity
can be reconstructed as the totality of preserved
determinations.

The role of determination in identity preservation
requires an independent ontological account.

\subsection{Relational Conditions of Identity Preservation}

An existent may be situated within a configuration
of relations that contributes to the conditions
of its existence.

Some relations may be constitutive of the particular
existent under examination. Others may change
without necessarily affecting its numerical identity.

The persistence of a relational configuration
may therefore be relevant to identity preservation,
but it cannot be treated as a universal criterion.

The same relational organization may be instantiated
by distinct existents. Conversely, an existent
may undergo relational transformation without
necessarily becoming numerically different.

The ontological relevance of relations must be
established according to their specific role,
domain, and respect.

It is also necessary to distinguish the preservation
of a relation from the preservation of the identity
of the relata. A relation may persist even when
one or more of its relata change, and the continuity
of a relation does not independently establish
the numerical identity of the entities involved.

Relational conditions may contribute to an account
of identity preservation, but their significance
must not be assumed in advance.

\subsection{Invariance as a Possible Condition of Identity Preservation}

Invariance may provide a relevant condition in
the investigation of identity preservation.

A feature that remains invariant under a specified
class of transformations may indicate that certain
ontological conditions have been preserved.

However, invariance is not identical with identity
preservation.

An invariant may be shared by distinct existents,
and the preservation of an invariant does not
independently establish numerical identity.
Conversely, the absence of an identified invariant
does not necessarily establish that identity has
ceased to persist.

The relevance of an invariant depends upon the
ontological role it plays in the existent under
examination.

An invariant may be necessary for identity preservation
within a specified domain if its constitutive
role has been established. It may also contribute
to a sufficient set of conditions when considered
alongside other ontological relations.

Nevertheless, no invariant should be treated as
universally necessary or sufficient without an
independent argument.

The distinction between formal invariance and
ontological invariance must also be maintained.
A formally preserved relation cannot automatically
be regarded as an ontologically preserved condition
of identity.

\subsection{Manifestation and the Conditions of Identity Preservation}

Manifestations may provide relevant indications
of the conditions under which an existent persists
through transformation.

The continuity or alteration of manifestations
may reveal changes in the determinations and
relations through which an existent is expressed.

However, manifestation continuity cannot be treated
as a sufficient condition of identity preservation.

A sequence of similar manifestations may belong
to numerically distinct existents. Conversely,
the manifestations associated with one existent
may undergo substantial alteration without
independently establishing the termination of
its identity.

The continuity of manifestation must therefore
be distinguished from the preservation of identity.

The same distinction applies to the capacity
to manifest. The preservation of such a capacity
may be relevant to the investigation, but it
does not independently establish numerical identity.

The conditions of identity preservation must
therefore be examined without identifying identity
with the continuity of its manifestations or
the uninterrupted exercise of its capacities.

\subsection{Causal Relations and Identity Preservation}

Causal relations may contribute to an account
of identity preservation by establishing connections
between successive conditions of an existent.

A sequence of transformations may be causally
connected, and such connections may help explain
how one condition gives rise to another.

Nevertheless, causal continuity does not independently
establish identity preservation.

Causal relations may connect numerically distinct
existents. A later existent may arise through
a causal process involving an earlier existent
without thereby being numerically identical with it.

Likewise, the absence of an established causal
connection does not independently settle the
question of identity discontinuity.

Causal relations must therefore be distinguished
from numerical identity.

Their ontological relevance depends upon whether
the causal relation is constitutive of the identity
under examination or merely describes a relation
between successive conditions.

\subsection{The Problem of Necessary Conditions}

A condition is necessary for identity preservation
within a specified domain when its absence excludes
the possibility of preserving the relevant identity
under the conditions being considered.

The identification of such conditions requires
an independent ontological argument.

A feature that frequently accompanies identity
preservation is not necessarily required for it.
A condition that is important in one class of
existents may be irrelevant or differently relevant
in another.

Consequently, no material, structural, functional,
relational, or manifestational condition should
be declared universally necessary merely because
it is associated with familiar cases of persistence.

The investigation must establish:

\begin{enumerate}
    \item the existent whose identity is under examination;
    \item the domain and transformations involved;
    \item the proposed necessary condition;
    \item the ontological reason why its absence would
    exclude identity preservation.
\end{enumerate}

Without these specifications, claims of necessity
remain insufficiently determined.

\subsection{The Problem of Sufficient Conditions}

A condition or set of conditions is sufficient
for identity preservation within a specified domain
when its presence provides adequate grounds for
establishing that the numerical identity of the
existent has remained the same.

Sufficiency must not be confused with correlation,
similarity, or the preservation of a frequently
observed feature.

A condition may be strongly associated with identity
preservation without being sufficient to establish it.
Even a combination of preserved features may fail
to distinguish one numerically identical existent
from another existent exhibiting the same configuration.

The identification of sufficient conditions therefore
requires an account of why the relevant conditions
exclude the possibility of numerical replacement
or the emergence of a distinct identity.

Such an account must be specific to the ontological
domain under examination.

It is possible that identity preservation depends
upon a conjunction of conditions rather than a
single feature. Yet this possibility must itself
be established and must not be treated as a
universal assumption.

\subsection{Conjunctive and Relational Conditions}

The conditions of identity preservation may be
conjunctive or relational.

A particular condition may be insufficient when
considered independently but may contribute to
a sufficient account when combined with other
conditions.

For example, material constitution, structural
organization, and relational configuration may
have different ontological roles. Their combined
preservation may be relevant in a particular
domain even when none of them independently
establishes identity preservation.

However, the mere accumulation of preserved
features does not automatically produce a sufficient
criterion.

A conjunction of conditions must explain why
their combined presence is ontologically decisive
for the identity under examination.

The relation among conditions must therefore
be established rather than inferred from their
number or descriptive completeness.

The investigation must also remain open to the
possibility that some conditions are dependent
upon others, that some are relevant only under
specific transformations, and that some may
have different significance across distinct
domains of existence.

\subsection{Identity Preservation Without Complete Feature Preservation}

Identity preservation must not be assumed to
require the preservation of every feature of
an existent.

An existent may undergo transformations in
its determinations, material constitution,
structural organization, relations, or manifestations.
The alteration of such features does not
automatically establish that its identity has
ceased.

The possibility of identity preservation through
incomplete feature preservation is therefore
ontologically significant.

However, this possibility must not be interpreted
as an unrestricted claim that identity persists
regardless of the transformations involved.

The absence of complete feature preservation
does not prove identity continuity, just as
the preservation of many features does not
independently establish it.

The question remains whether the conditions
that have changed are compatible with the
persistence of the identity under examination.

This requires distinguishing between transformable
determinations and conditions whose alteration
would be incompatible with the continued identity
of the existent.

The distinction cannot be settled through
a universal list of preserved or transformable
features. It must be established in relation
to the ontological constitution of the existent
and the transformations involved.

\subsection{Identity Preservation Across Interruption}

Identity preservation may also be examined
under conditions of interruption.

An existent may undergo an interruption in
a particular process, manifestation, relation,
or mode of activity. Such interruption may
alter the conditions under which persistence
is described.

Nevertheless, the interruption of a particular
process or manifestation does not independently
establish the termination of identity.

Nor does the possibility of resumption establish
that the same identity has persisted throughout
the interruption.

The ontological relation between the conditions
before, during, and after interruption must
be examined independently.

The conditions of identity preservation must
therefore account for interruption without
equating it either with identity termination
or with an insignificant temporary absence.

The relevance of interruption depends upon
the particular existent, the interrupted
condition, and the respect under which identity
preservation is being examined.

\subsection{The Possibility of Multiple Identity Relations}

The conditions of identity preservation must
not presuppose that every existent is associated
with exactly one identity.

The relation between an existent and its identity
or identities remains an ontological question.
The preservation of one identity relation does
not automatically settle whether other identity
relations may also be involved.

Similarly, the multiplicity of components,
determinations, manifestations, or relations
does not independently establish a plurality
of identities.

The conditions of identity preservation must
therefore remain neutral concerning the number
of identities until the relevant ontological
relations have been established.

This neutrality is necessary to prevent a
particular model of unity or plurality from
determining the conclusions in advance.

\subsection{The Subject, Domain, and Respect of Identity Preservation}

Every claim concerning identity preservation
must specify its subject, domain, and respect.

The subject identifies the existent or identity
whose preservation is being examined.

The domain identifies the relevant class of
transformations, conditions, or temporal
succession.

The respect identifies the ontological aspect
under which preservation is asserted or denied.

A condition may be necessary in one domain
without being necessary in another. A feature
may be sufficient under a restricted class
of transformations without providing a universal
criterion.

Likewise, the preservation of identity in one
respect must not be generalized into a claim
of complete preservation in every respect.

These specifications are not merely methodological
conveniences. They are necessary for determining
the ontological scope and validity of claims
concerning identity preservation.

\subsection{Methodological Requirements}

The investigation of the conditions of identity
preservation must observe the following requirements:

\begin{enumerate}
    \item Identity preservation must be distinguished
    from the preservation of individual features.
    \item No material, structural, functional,
    relational, or manifestational condition
    should be assumed universally necessary
    or sufficient.
    \item The ontological relevance of every
    proposed condition must be independently
    established.
    \item Necessary conditions must be distinguished
    from sufficient conditions and from conditions
    that merely accompany persistence.
    \item Conjunctive conditions must be justified
    through their ontological relations rather
    than their descriptive accumulation.
    \item Invariance and manifestation continuity
    must not be treated as automatic criteria
    of identity preservation.
    \item Identity preservation through transformation
    must remain distinct from unrestricted persistence.
    \item The possibility of interruption must be
    examined without presupposing identity termination
    or uninterrupted preservation.
    \item The number of identities involved must
    remain open to independent investigation.
    \item Every claim must specify its subject,
    domain, and respect.
\end{enumerate}

These requirements establish a framework for
investigating identity preservation without
reducing it to any single form of continuity
or feature preservation.

\subsection{Scope of the Present Section}

This section has examined the ontological conditions
under which identity preservation may be established.

Material constitution, structural organization,
determination, relational configuration, invariance,
manifestation, and causal connection may contribute
to an account of identity preservation. Nevertheless,
none of these conditions can be treated as a
universal criterion without an independent
ontological justification.

The distinction between necessary and sufficient
conditions is essential. A feature that accompanies
identity preservation is not necessarily required
for it, and a feature that is required in a
particular domain may not be sufficient to establish
identity preservation.

The preservation of identity must therefore be
understood through an account of the ontological
relations that make the continued numerical identity
of a particular existent possible under specified
transformations.

The next section examines the limits of identity
continuity. It investigates the boundaries of
what can be established about identity preservation,
the conditions under which continuity claims remain
underdetermined, and the limits of extending
particular criteria across different domains
of existence.

% ============================================================
% End of Section 05
% ============================================================